\NSPage{139}%
The heat profile approaches the original power law at large \NSi{NS-F-P139-001}{X}, with the derivative bounds
in \eqref{eq:A.38}. Denote by \NSi{NS-F-P139-002}{E_{\mathrm{cl}}} the reference profile constructed in
Proposition~\ref{prop:A.4}. We use this small change to replace the terminal profile and restore the three affected
moments by earlier corrections. The other two moments are unchanged because both modified regions have
\NSi{NS-F-P139-003}{U=0}.

\begin{proposition}\label{prop:A.7}
Let \NSi{NS-F-P139-004}{X_K=e^{.2}X_{\mathrm{tail}}}, and choose a smooth step
\NSi{NS-F-P139-005}{0\leq\chi_K(y)\leq 1} equal to zero for
\NSi{NS-F-P139-006}{y\leq .2} and to one for \NSi{NS-F-P139-007}{y\geq .5}. The terminal replacement
\NSFormula{NS-F-P139-008}{display}{A.39}
\begin{equation}
  E_{\mathrm{cl}}\longmapsto E_{\mathrm{cl}}
  \bigl[1+\chi_K(y)(\mathcal H(2d/X)-1)\bigr]
  \tag{A.39}\label{eq:A.39}
\end{equation}
can be compensated by three additive \NSi{NS-F-P139-009}{E}-bumps in the second reserved patch in the intermediate power-law
interval. The complete edit preserves the functions \NSi{NS-F-P139-010}{M,J}, the totals
\NSi{NS-F-P139-011}{S(\infty),C_p(\infty)} from \eqref{eq:4.15}, and the renormalized angular moment
\NSi{NS-F-P139-012}{\int_0^\infty(H-H_{\mathrm{pow}})\,dX}, pointwise on
\NSi{NS-F-P139-013}{[-1,1]}. For sufficiently large \NSi{NS-F-P139-014}{X_R}, it preserves
\NSi{NS-F-P139-015}{E>0} and the strict admissible stress cone from that patch through
\NSi{NS-F-P139-016}{y=.5}. The original axis pressure datum and the profiles before the compensation patch are unchanged.
\end{proposition}

\begin{proof}
We estimate the three changes caused by the heat factor, cancel them with three earlier angular bumps, and use the
normalized radial moment formulas to preserve the admissible stress cone.

Put \NSi{NS-F-P139-017}{x=X/X_K} and let \NSi{NS-F-P139-018}{e_K} be the reference amplitude at
\NSi{NS-F-P139-019}{X_K}. The reference tail is comparable to
\NSi{NS-F-P139-020}{e_Kx^{-A}}. Equations~\eqref{eq:A.34} to~\eqref{eq:A.36} and the chain rule show that,
for all fixed integers \NSi{NS-F-P139-021}{j,m\geq 0},
\NSFormula{NS-F-P139-022}{display}{none}
\[
  \bigl|(X\partial_X)^j\partial_\eta^m(E-E_{\mathrm{cl}})\bigr|
  \leq C_{j,m}e_KX_K^{-1}x^{-A-1},\qquad x\geq 1.
\]
The constants may depend on the already fixed schedule. In particular, all derivatives in this estimate are bounded at
\NSi{NS-F-P139-023}{d=0}. Writing \NSi{NS-F-P139-024}{\Delta} for the change caused just by \eqref{eq:A.39}, the three
discrepancies satisfy
\NSFormula{NS-F-P139-025}{display}{A.40}
\begin{equation}
  \left|\partial_\eta^m\Delta\int_0^\infty\frac{E^2}{2X}\,dX\right|
  \leq C_m\frac{e_K^2}{X_K}\int_1^\infty x^{-2A-2}\,dx,
  \tag{A.40}\label{eq:A.40}
\end{equation}
\NSFormula{NS-F-P139-026}{display}{A.41}
\begin{equation}
  \left|\partial_\eta^m\Delta S(\infty)\right|
  \leq C_me_K^2\int_1^\infty x^{-2A-1}\,dx,
  \tag{A.41}\label{eq:A.41}
\end{equation}
\NSFormula{NS-F-P139-027}{display}{A.42}
\begin{equation}
  \left|\partial_\eta^m\Delta\int_0^\infty(H-H_{\mathrm{pow}})\,dX\right|
  \leq C_me_KX_K^{1/2}\int_1^\infty x^{-A-1/2}\,dx
  \leq C_{m,h}e_KX_K^{1/2}.
  \tag{A.42}\label{eq:A.42}
\end{equation}
Here \NSi{NS-F-P139-028}{H=\sqrt{2X}E} and
\NSi{NS-F-P139-029}{H_{\mathrm{pow}}=\sqrt{2X}E_{\mathrm{pow}}}; the pure power is unchanged. The final integral is
\NSi{NS-F-P139-030}{1/h}. It is finite because \NSi{NS-F-P139-031}{h>0} has already been fixed before the large-radius choice.

The three moment discrepancies are small after normalization by their natural radial scales. We cancel them on the second
reserved patch, where the power-law profile gives three independent moment weights. At the start
\NSi{NS-F-P139-032}{X_*} of this patch, put \NSi{NS-F-P139-033}{x_*=X/X_*}. Its unmodified profile is
\NSi{NS-F-P139-034}{E=e_*f(\eta)x_*^{-1/2-\lambda}}, \NSi{NS-F-P139-035}{U=0}, where
\NSi{NS-F-P139-036}{f=(1+\eta^2)^{-1}\geq 1/2}. Choose three nonnegative bumps
\NSi{NS-F-P139-037}{\beta_i} with ordered disjoint supports inside this patch and add
\NSFormula{NS-F-P139-038}{display}{none}
\[
  \delta E=e_*\sum_{i=1}^3c_i(\eta)\beta_i(x_*).
\]
Normalize the pressure, \NSi{NS-F-P139-039}{S}, and angular rows respectively by
\NSFormula{NS-F-P139-040}{display}{A.43}
\begin{equation}
  e_*^2,\qquad X_*e_*^2,\qquad X_*^{3/2}e_*.
  \tag{A.43}\label{eq:A.43}
\end{equation}
Their derivatives at \NSi{NS-F-P139-041}{c=0} are the integrals against the weights
\NSFormula{NS-F-P139-042}{display}{none}
\[
  fx_*^{-3/2-\lambda},\qquad -fx_*^{-1/2-\lambda},\qquad \sqrt{2}x_*^{1/2}.
\]

\NSPage{140}%
The exponents are distinct, so Lemma~\ref{lem:A.1} makes this matrix invertible. Its inverse is uniform in
\NSi{NS-F-P140-001}{\eta\in[-1,1]}. The first two rows have quadratic remainders and the last is linear. Since
\NSi{NS-F-P140-002}{X_*/X_K} and \NSi{NS-F-P140-003}{e_*/e_K} are fixed, dividing
Equations~\eqref{eq:A.40} to~\eqref{eq:A.42} by \eqref{eq:A.43} gives
\NSi{NS-F-P140-004}{O_m(X_K^{-1})} after every fixed number of \NSi{NS-F-P140-005}{\eta}-derivatives. The local existence
result of Lemma~\ref{lem:A.2} gives a smooth solution with
\NSi{NS-F-P140-006}{\lVert\partial_\eta^m c\rVert_\infty\leq C_mX_K^{-1}}. Choose
\NSi{NS-F-P140-007}{X_R} sufficiently large that the moment equations have a unique solution near zero correction
coefficients, with invertible Jacobian. Implicit differentiation then gives a finite bound at each fixed higher order of
\NSi{NS-F-P140-008}{\eta}-derivatives. Only the finitely many derivative bounds used to preserve the cone must meet prescribed
smallness thresholds.

The three altered moments have now been restored, and \NSi{NS-F-P140-009}{M,J} remain unchanged because both edited regions
have \NSi{NS-F-P140-010}{U=0}. We must still check that the local changes in the profiles and cumulative integrals preserve the
cone between the compensation patch and the terminal interval. The normalized radial moment formulas \eqref{eq:4.16}, in
\NSi{NS-F-P140-011}{X/X_R}, transfer the small profile and moment changes to \NSi{NS-F-P140-012}{Q_s,N_s} on the fixed compact
log interval from this patch to \NSi{NS-F-P140-013}{y=.5}. More explicitly,
\NSi{NS-F-P140-014}{I=X_R^{3/2}\widehat I},
\NSi{NS-F-P140-015}{J=X_R^{3/2}\widehat J},
\NSi{NS-F-P140-016}{M=X_R\widehat M}, and
\NSi{NS-F-P140-017}{S=X_R\widehat S}; insertion in the identities for
\NSi{NS-F-P140-018}{Q_s,N_s} cancels every \NSi{NS-F-P140-019}{X_R}-factor, including those in
\NSi{NS-F-P140-020}{H}. The comparison costs one further \NSi{NS-F-P140-021}{\eta}-derivative and gives
\NSi{NS-F-P140-022}{O(X_K^{-1})} changes in the required \NSi{NS-F-P140-023}{\eta}-derivatives of
\NSi{NS-F-P140-024}{Q_s,N_s}. The positive lower bounds for \NSi{NS-F-P140-025}{E,Q_s}, the slope bounds, and the admissible
stress cone margins of Proposition~\ref{prop:A.4} consequently persist for large \NSi{NS-F-P140-026}{X_R}. The intermediate
values of \NSi{NS-F-P140-027}{I} and the transition value of \NSi{NS-F-P140-028}{Q_s} may vary while these admissible stress cone
inequalities persist.

Finally, zero total pressure discrepancy means that integration of
\NSi{NS-F-P140-029}{\Pi_X=E^2/(2X)} forward from the unchanged \NSi{NS-F-P140-030}{\Pi_0} still gives the pressure normalized to
vanish at radial infinity. Before the compensation patch this forward solution is unchanged. In particular the terminal heat
factor, which need only be smooth in \NSi{NS-F-P140-031}{\eta}, does not alter the analytic axis datum.
\end{proof}

\subsection{Recovering the stress from the exterior.}\label{sec:A.7}
In this subsection we will use the corrected moments to prove that the stress vanishes beyond the annulus
(Lemma~\ref{lem:A.8}) and to determine its direction as the outer edge is approached
(Proposition~\ref{prop:A.10}). We will first show that the moment identities make the total weighted integrals of the leading
residual zero. This will allow us to compute the stress by integration from the current radius to infinity. This representation
depends only on the exterior profile once the regular axis and matching moment conditions have been supplied by
Proposition~\ref{prop:B.2} and Corollary~\ref{cor:B.10}.

\begin{lemma}\label{lem:A.8}
Suppose the leading axisymmetric field is smooth at the axis, has \NSi{NS-F-P140-032}{U=0} and
\NSi{NS-F-P140-033}{E=E_{\mathrm{pow}}f_o[1+\chi_K(\mathcal H(2d/X)-1)]} for
\NSi{NS-F-P140-034}{y\geq 0}, and satisfies the exact moment conditions
\NSFormula{NS-F-P140-035}{display}{none}
\[
  M(\infty)=J(\infty)=S(\infty)=0,\qquad
  \int_0^\infty(H-H_{\mathrm{pow}})\,dX=0,\qquad
  \Pi(X)=-\frac12\int_X^\infty\frac{E(x)^2}{x}\,dx.
\]
The leading tangential stresses vanish for \NSi{NS-F-P140-036}{X\geq X_b}. On
\NSi{NS-F-P140-037}{y\geq 1/2}, where \NSi{NS-F-P140-038}{\chi_K=1}, they are given by the weighted integrals of the leading
residual from \NSi{NS-F-P140-039}{r} to infinity, with positive integration sign, in \eqref{eq:A.46}.
\end{lemma}

\begin{proof}
The radial moment functions in the tail admit the following representations by integrals from \NSi{NS-F-P140-040}{X} to
infinity. Since \NSi{NS-F-P140-041}{U=M=J=0} there, \NSi{NS-F-P140-042}{W=1}, and
\NSFormula{NS-F-P140-043}{display}{none}
\[
  I(X)=\frac{XH_{\mathrm{pow}}(X)}{1-h}-\int_X^\infty(H-H_{\mathrm{pow}})\,dx,
  \qquad
  S(X)=\frac12\int_X^\infty E^2\,dx,
\]
\NSFormula{NS-F-P140-044}{display}{none}
\[
  Q_s=-1+\frac{(1-h)I-D\eta I_\eta}{XH},
  \qquad
  N_s=\frac{4h\eta S-dS_\eta}{X}+4A\eta\Pi-d\Pi_\eta.
\]

\NSPage{141}%
The first identity uses \NSi{NS-F-P141-001}{h<1} to integrate \NSi{NS-F-P141-002}{H_{\mathrm{pow}}} from zero. The others
follow directly from the exact moments and the integral identities for \NSi{NS-F-P141-003}{Q_s,N_s} \eqref{eq:4.16}.

We verify that no boundary term obstructs the backward stress formula. Let
\NSi{NS-F-P141-004}{K_{\mathrm{pow}}=c_\infty s^{-A}}. At large radius, uniformly for \NSi{NS-F-P141-005}{\tau} in bounded
nonnegative intervals,
\NSFormula{NS-F-P141-006}{display}{A.44}
\begin{equation}
  K-K_{\mathrm{pow}}=O_h(\tau r^{-3-2h}),\qquad
  \partial_tK=O_h(r^{-3-2h}),\qquad
  r^2K_r-rK=O_h(r^{-2h}).
  \tag{A.44}\label{eq:A.44}
\end{equation}
These follow from \eqref{eq:A.36} and its fixed derivative versions. In particular, at any sufficiently large fixed physical
radius \NSi{NS-F-P141-007}{R_*},
\NSFormula{NS-F-P141-008}{display}{none}
\[
  \int_{R_*}^\infty\bigl(r^2|K-K_{\mathrm{pow}}|+|\partial_tK|\bigr)\,dr
  \leq C_h(1+\tau)\int_{R_*}^\infty r^{-1-2h}\,dr
  =\frac{C_h(1+\tau)}{2h}R_*^{-2h}.
\]
Thus the subtracted angular moment is a convergent integral. Its time derivative is justified by the same integrable majorant,
uniformly on compact time intervals. We obtain
\NSFormula{NS-F-P141-009}{display}{A.45}
\begin{equation}
  \int_0^\infty r^2(u_\theta-K_{\mathrm{pow}})\,dr
  =q^{3/2-A}\int_0^\infty(H-H_{\mathrm{pow}})\,dX=0.
  \tag{A.45}\label{eq:A.45}
\end{equation}
At the axis the subtracted power has weight \NSi{NS-F-P141-010}{r^{1-2h}}, also integrable. Its time derivative is zero. The
angular transport integral is the sum of the radial boundary
\NSi{NS-F-P141-011}{[r^2u_ru_\theta]_0^\infty} and the axial derivative of
\NSi{NS-F-P141-012}{\int r^2u_zu_\theta\,dr=q^{3/2-2A}J(\infty)=0}. The radial boundary vanishes by axis regularity and by
\NSi{NS-F-P141-013}{U=M=0} on the exterior, which also gives \NSi{NS-F-P141-014}{V_0^{(0)}=0}. The radial viscosity integrates
to \NSi{NS-F-P141-015}{[r^2\partial_ru_\theta-ru_\theta]_0^\infty=0} by \eqref{eq:A.44}. Thus the leading angular residual
\NSi{NS-F-P141-016}{\mathscr R_\theta^{(0)}} has zero integral against \NSi{NS-F-P141-017}{r^2\,dr}.

For the axial component, \NSi{NS-F-P141-018}{\int ru_z\,dr=q^{1-A}M(\infty)=0}. The pressure identity
\NSi{NS-F-P141-019}{p(r)=-\int_r^\infty u_\theta(r')^2\,dr'/r'} gives
\NSFormula{NS-F-P141-020}{display}{none}
\[
  \int_0^\infty r(u_z^2+p)\,dr
  =q^{1-2A}\int_0^\infty(U^2-E^2/2)\,dX=0.
\]
Indeed the exact pressure calculation is
\NSFormula{NS-F-P141-021}{display}{none}
\[
  \int_0^\infty rp(r)\,dr
  =-\int_0^\infty\frac{u_\theta(r')^2}{r'}
    \left(\int_0^{r'}r\,dr\right)dr'
  =-\frac12\int_0^\infty r'u_\theta(r')^2\,dr'.
\]
The exterior decay \NSi{NS-F-P141-022}{p=O(r^{-2-4h})} ensures convergence and
\NSi{NS-F-P141-023}{r^2p\to 0}. The boundary terms from radial transport and viscosity in the axial equation also vanish.
Hence \NSi{NS-F-P141-024}{\int r\mathscr R_z^{(0)}\,dr=0}.

Write \NSi{NS-F-P141-025}{T_\theta,T_z} for the physical stress components, so
\NSi{NS-F-P141-026}{T=q^{-A-1/2}\mathcal T_0}. Their defining integrals from the axis can now be rewritten as
\NSFormula{NS-F-P141-027}{display}{A.46}
\begin{equation}
  T_\theta(r)=\frac1{r^2}\int_r^\infty r'^2\mathscr R_\theta^{(0)}(r')\,dr',
  \qquad
  T_z(r)=\frac1r\int_r^\infty r'\mathscr R_z^{(0)}(r')\,dr'.
  \tag{A.46}\label{eq:A.46}
\end{equation}
For \NSi{NS-F-P141-028}{X\geq X_b}, the field is
\NSi{NS-F-P141-029}{(u_r,u_\theta,u_z)=(0,K,0)} with \NSi{NS-F-P141-030}{z}-independent pressure normalized to vanish at
radial infinity. Lemma~\ref{lem:A.6} makes both leading residuals zero there. Thus these stress components vanish for
\NSi{NS-F-P141-031}{X\geq X_b}.
\end{proof}

The backward integrals contain the cutoff factor that vanishes to every order at the outer edge. The next lemma separates this
factor from a smooth coefficient, so that the stress direction can be estimated even where the stress itself tends to zero.

\begin{lemma}\label{lem:A.9}
Let \NSi{NS-F-P141-032}{c>0}, let \NSi{NS-F-P141-033}{j} be a fixed nonnegative integer, and let
\NSi{NS-F-P141-034}{b(u,\eta)} be smooth for \NSi{NS-F-P141-035}{0\leq u\leq\delta_0} and for
\NSi{NS-F-P141-036}{\eta} in a compact interval \NSi{NS-F-P141-037}{K}. There exists a coefficient
\NSi{NS-F-P141-038}{B} such that, for \NSi{NS-F-P141-039}{0<\delta\leq\delta_0},
\NSFormula{NS-F-P141-040}{display}{A.47}
\begin{equation}
  \int_0^\delta e^{-c/u^2}u^{-j}b(u,\eta)\,du
  =\frac12e^{-c/\delta^2}\delta^{3-j}B(\delta,\eta).
  \tag{A.47}\label{eq:A.47}
\end{equation}

\NSPage{142}%
The coefficient \NSi{NS-F-P142-001}{B} is smooth on
\NSi{NS-F-P142-002}{[0,\delta_0]\times K}, has
\NSi{NS-F-P142-003}{B(0,\eta)=b(0,\eta)/c}, and has bounded mixed derivatives of every fixed order there. The same assertion
holds with any finite collection of smooth compact parameters in place of \NSi{NS-F-P142-004}{\eta}.
\end{lemma}

\begin{proof}
Substitute \NSi{NS-F-P142-005}{u=\delta/(1+\delta^2v)^{1/2}}. The Jacobian is
\NSi{NS-F-P142-006}{-\tfrac12\delta^3(1+\delta^2v)^{-3/2}}, and
\NSi{NS-F-P142-007}{c/u^2=c/\delta^2+cv}. This proves \eqref{eq:A.47} with
\NSFormula{NS-F-P142-008}{display}{none}
\[
  B(\delta,\eta)=\int_0^\infty e^{-cv}(1+\delta^2v)^{(j-3)/2}
  b\!\left(\frac{\delta}{\sqrt{1+\delta^2v}},\eta\right)\,dv.
\]
Every fixed mixed derivative of the final integrand is bounded by \NSi{NS-F-P142-009}{e^{-cv}} times a polynomial in
\NSi{NS-F-P142-010}{v}, uniformly on the closed parameter rectangle. Dominated differentiation proves smoothness and the
derivative bounds. Setting \NSi{NS-F-P142-011}{\delta=0} under the integral gives the stated endpoint value.
\end{proof}

The stress vanishes at the outer edge, so its limiting direction depends on the relative rates at which its two components
vanish. The backward formula and the flat integral lemma give these rates, with a positive leading coefficient in the angular
component.

\begin{proposition}\label{prop:A.10}
Under the regular axis and exact moment hypotheses of Lemma~\ref{lem:A.8}, the profile on the terminal interval satisfies the
admissible stress cone on \NSi{NS-F-P142-012}{.5\leq y<3}. Its stress direction extends smoothly to the outer endpoint, and
\NSi{NS-F-P142-013}{2-(a-2)(T_z/T_\theta)^2}, interpreted through this extension, has a uniform positive lower bound. More
precisely, for \NSi{NS-F-P142-014}{\delta=3-y} sufficiently small,
\NSFormula{NS-F-P142-015}{display}{A.48}
\begin{equation}
  \mathcal T_{0,\theta}=e^{-4/\delta^2}\delta^{-3}b_\theta(\delta,\eta),
  \qquad b_\theta(0,\eta)>0,
  \tag{A.48}\label{eq:A.48}
\end{equation}
\NSFormula{NS-F-P142-016}{display}{A.49}
\begin{equation}
  \mathcal T_{0,z}=e^{-4/\delta^2}\delta^3b_z(\delta,\eta),
  \tag{A.49}\label{eq:A.49}
\end{equation}
\NSFormula{NS-F-P142-017}{display}{A.50}
\begin{equation}
  \frac{\mathcal T_{0,z}}{\mathcal T_{0,\theta}}=\delta^6\frac{b_z}{b_\theta}\longrightarrow 0.
  \tag{A.50}\label{eq:A.50}
\end{equation}
The coefficients are smooth on a closed outer collar, and \NSi{NS-F-P142-018}{b_\theta} has a positive lower bound there. For
every fixed mixed profile derivative \NSi{NS-F-P142-019}{\partial^I}, there are finite constants
\NSi{NS-F-P142-020}{C_I,N_I} such that
\NSFormula{NS-F-P142-021}{display}{A.51}
\begin{equation}
  |\partial^I\mathcal T_0|\leq C_Ie^{-4/\delta^2}\delta^{-N_I},
  \qquad |\mathcal T_0|\geq ce^{-4/\delta^2}\delta^{-3}.
  \tag{A.51}\label{eq:A.51}
\end{equation}
These constants may depend on all fixed profile choices and on \NSi{NS-F-P142-022}{I}, but not on the physical scale
\NSi{NS-F-P142-023}{q}.
\end{proposition}

\begin{proof}
We first derive a positive angular stress and bound the ratio of axial to angular stress. We then factor the flat endpoint terms
to prove smoothness of the limiting direction.

On \NSi{NS-F-P142-024}{y\geq 1/2}, where \NSi{NS-F-P142-025}{\chi_K=1}, let
\NSi{NS-F-P142-026}{f=f_o(y)} and write \NSi{NS-F-P142-027}{f'=\partial_yf}. Here
\NSi{NS-F-P142-028}{u_\theta=Kf} and \NSi{NS-F-P142-029}{u_r=u_z=0}. The heat equation from Lemma~\ref{lem:A.6} and
\NSi{NS-F-P142-030}{y_t=(qL)^{-1}},
\NSi{NS-F-P142-031}{y_z=-2\eta/(q^DL)} yield
\NSFormula{NS-F-P142-032}{display}{A.52}
\begin{equation}
  \mathscr R_\theta^{(0)}
  =\frac{Kf'}{qL}-K(f_{rr}+r^{-1}f_r)-2K_rf_r,
  \tag{A.52}\label{eq:A.52}
\end{equation}
\NSFormula{NS-F-P142-033}{display}{A.53}
\begin{equation}
  p_z(r)=\frac{2\eta}{q^DL}\int_y^3K(y')^2f(y')f'(y')\,dy'.
  \tag{A.53}\label{eq:A.53}
\end{equation}
Integrating the viscous terms of \eqref{eq:A.52} by parts in \eqref{eq:A.46} gives the exact formula
\NSFormula{NS-F-P142-034}{display}{A.54}
\begin{equation}
  T_\theta
  =Kf_r+\frac1{r^2}\int_r^\infty(r'K-r'^2K_{r'})f_{r'}\,dr'
  +\frac1{qLr^2}\int_r^\infty r'^2Kf'\,dr'.
  \tag{A.54}\label{eq:A.54}
\end{equation}
All three terms are nonnegative: \NSi{NS-F-P142-035}{K>0}, \NSi{NS-F-P142-036}{K_r<0}, and
\NSi{NS-F-P142-037}{f'\geq 0}. On the remaining log interval, whose length is at most
\NSi{NS-F-P142-038}{2.5}, the ratios of radii and heat factors are bounded above and below.

\NSPage{143}%
Since \NSi{NS-F-P143-001}{\int_y^3f'(v)\,dv=\rho_o\psi_o(y)}, the last term gives
\NSFormula{NS-F-P143-002}{display}{none}
\[
  T_\theta\geq c\frac{rK}{qL}\rho_o\psi_o(y)>0\qquad(y<3).
\]
Similarly \eqref{eq:A.53} and the integral representation of \NSi{NS-F-P143-003}{T_z} imply
\NSFormula{NS-F-P143-004}{display}{A.55}
\begin{equation}
  |p_z|\leq Cq^{-D}K^2\rho_o\psi_o,
  \qquad |T_z|\leq Crq^{-D}K^2\rho_o\psi_o,
  \qquad \left|\frac{T_z}{T_\theta}\right|\leq Cq^{1-D}K
  =CE_{\mathrm{pow}}\mathcal H(2d/X).
  \tag{A.55}\label{eq:A.55}
\end{equation}
The constants here are uniform for sufficiently large \NSi{NS-F-P143-005}{X_R} and small \NSi{NS-F-P143-006}{h}. The last
equality uses \NSi{NS-F-P143-007}{1-D=A}; the amplitude reduction in the exterior transition of \eqref{eq:A.17} makes its last
expression small. Meanwhile \NSi{NS-F-P143-008}{b_s=0}, and
\NSFormula{NS-F-P143-009}{display}{A.56}
\begin{equation}
  a=2+2h+2Z\frac{\mathcal H'}{\mathcal H}-2\frac{f'}f>2+h,
  \qquad Z=2d/X,
  \tag{A.56}\label{eq:A.56}
\end{equation}
because \NSi{NS-F-P143-010}{f'/f<h/4} and, for large \NSi{NS-F-P143-011}{X_R},
\NSi{NS-F-P143-012}{-Z\mathcal H'/\mathcal H<h/4}. Also \NSi{NS-F-P143-013}{a\leq 2+2h}. Thus
\NSi{NS-F-P143-014}{v_s=a},
\NSi{NS-F-P143-015}{P_{\mathrm c}-v_s=\mathcal T_{0,\theta}/F>0}, and
\NSi{NS-F-P143-016}{J_{\mathrm c}=\mathcal T_{0,z}/F}. The directional cone inequality is
\NSi{NS-F-P143-017}{(a-2)(T_z/T_\theta)^2<2}, with fixed strict slack by \eqref{eq:A.55}. This proves the admissible stress
cone at each point before the outer endpoint \NSi{NS-F-P143-018}{X=X_b}.

The cone inequality now holds at every interior point of the terminal interval. It remains to prove that the stress direction
extends smoothly to the endpoint. We factor the common vanishing part of the two stress components and show that the angular
coefficient stays positive. The specified smooth step gives, near \NSi{NS-F-P143-019}{\delta=0},
\NSFormula{NS-F-P143-020}{display}{none}
\[
  \psi_o=e^{-4/\delta^2}g(\delta),
  \qquad
  g(\delta)=\frac{1}{e^{-(1-\delta/2)^{-2}}+e^{-4/\delta^2}},
  \qquad g(0)=e,
\]
and hence
\NSFormula{NS-F-P143-021}{display}{none}
\[
  f'=\rho_oe^{-4/\delta^2}\delta^{-3}\bigl[8g(\delta)+\delta^3g'(\delta)\bigr].
\]
We apply Lemma~\ref{lem:A.9} to the backward stress formulas.

Multiplying the boundary term in \eqref{eq:A.54} by \NSi{NS-F-P143-022}{q^{A+1/2}} gives
\NSFormula{NS-F-P143-023}{display}{none}
\[
  q^{A+1/2}Kf_r=\frac{2E_{\mathrm{pow}}(X)\mathcal H(2d/X)}{\sqrt{2X}}f'.
\]
The two integral terms contain \NSi{NS-F-P143-024}{f'}, so \eqref{eq:A.47} with
\NSi{NS-F-P143-025}{j=3,c=4} gives a factor \NSi{NS-F-P143-026}{e^{-4/\delta^2}} times a smooth coefficient. They therefore
vanish to order at least \NSi{NS-F-P143-027}{\delta^3} relative to the boundary factor. It follows that \eqref{eq:A.48} holds,
with
\NSFormula{NS-F-P143-028}{display}{none}
\[
  b_\theta(0,\eta)=\frac{16\rho_oE_{\mathrm{pow}}(X_b)}{\sqrt{2X_b}}
  \mathcal H(2d/X_b)g(0)>0.
\]
The positive lower bound is uniform for \NSi{NS-F-P143-029}{-1\leq\eta\leq 1}, since \NSi{NS-F-P143-030}{\mathcal H} has a
positive minimum on \NSi{NS-F-P143-031}{[0,2/X_b]}.

The pressure integral \eqref{eq:A.53} has one \NSi{NS-F-P143-032}{f'}-factor, so the same identity first gives
\NSi{NS-F-P143-033}{p_z=e^{-4/\delta^2}} times a smooth coefficient, apart from its factor
\NSi{NS-F-P143-034}{q^{-D-2A}}. The integral representation of \NSi{NS-F-P143-035}{T_z} requires one further integration with
smooth weights. Applying \eqref{eq:A.47} with \NSi{NS-F-P143-036}{j=0} gives precisely \eqref{eq:A.49}. In this normalization
the powers cancel as
\NSi{NS-F-P143-037}{q^{A+1/2}q^{1/2-D-2A}=q^{1-D-A}=1}. All remaining coefficients contain only smooth functions of
\NSi{NS-F-P143-038}{X,\eta,\mathcal H(2d/X)}, and \NSi{NS-F-P143-039}{L^{-1}}; here
\NSi{NS-F-P143-040}{L\geq 1-2h>0}. This also verifies closed-parameter smoothness and the absence of a physical scale in both
flat factors.

\NSPage{144}%
Equation~\eqref{eq:A.50} follows by division by the positive coefficient \NSi{NS-F-P144-001}{b_\theta}. Its limiting direction
is \NSi{NS-F-P144-002}{(1,0)}, where the normalized directional quadratic expression equals \NSi{NS-F-P144-003}{2}. The
stress itself is zero at the endpoint. Finally each fixed profile derivative of the factorizations only introduces a finite
inverse power of \NSi{NS-F-P144-004}{\delta}, proving \eqref{eq:A.51}. A smooth positive factor from the other endpoint can
multiply this outer weight without changing these local conclusions.
\end{proof}

\section{Analytic profiles near the axis and their continuation}\label{sec:B}

In this appendix we construct the inner part of the leading profile \NSi{NS-F-P144-005}{(U,E,\Pi)}, where the leading residual
stress \NSi{NS-F-P144-006}{\mathcal T_0} must vanish (Proposition~\ref{prop:B.2}). The outer profile has already fixed the pressure datum
\NSi{NS-F-P144-007}{\Pi_0} at the axis (Lemma~\ref{lem:A.5}). With that datum, we choose the remaining axis data
\NSi{NS-F-P144-008}{U_*,\phi_*} and solve the nonlinear profile equations \eqref{eq:4.13} on a short radial interval. The solution
must be regular at the axis and have the strict shear inequality \eqref{eq:B.19} at its outer endpoint, so that it can be
continued toward the outer profile (Corollary~\ref{cor:B.10}). A large parameter \NSi{NS-F-P144-009}{\Lambda} in the azimuthal
datum \eqref{eq:B.3} gives both the analytic solution and the endpoint bounds needed for this continuation
(Propositions~\ref{prop:B.2} and~\ref{prop:B.3}). The corresponding velocity is smooth in Cartesian coordinates across the axis
by Equations~\eqref{eq:4.4} to~\eqref{eq:4.5}.

\subsection{Choice of axis data.}\label{sec:B.1}
The pressure datum \NSi{NS-F-P144-010}{\Pi_0} is prescribed, while the azimuthal datum
\NSi{NS-F-P144-011}{\phi_*=\phi(0,\eta)} and the axial velocity \NSi{NS-F-P144-012}{U_*} at the axis remain to be chosen. In this
subsection we choose these data to keep the angular source \NSi{NS-F-P144-013}{S_q} positive and enforce the endpoint inequality
throughout the parameter interval (Proposition~\ref{prop:B.3}). The azimuthal and axial terms will provide the required lower
bound on complementary parts of that interval, as separated in \eqref{eq:B.2}.

We use the datum \NSi{NS-F-P144-014}{\Pi_0} of Lemma~\ref{lem:A.5}; in particular it is real analytic near
\NSi{NS-F-P144-015}{[-1,1]}, is even, and satisfies
\NSFormula{NS-F-P144-016}{display}{none}
\[
  \Pi_0\leq-cP_*^2f^2,
  \qquad \eta\Pi_{0\eta}>0\quad(\eta\neq 0),
  \qquad f=(1+\eta^2)^{-1}.
\]
All schedule parameters, including \NSi{NS-F-P144-017}{0<h\leq 10^{-2}}, have already been fixed. Choose
\NSi{NS-F-P144-018}{0<j_0\leq .05} small and set
\NSFormula{NS-F-P144-019}{display}{B.1}
\begin{equation}
\begin{aligned}
  U_*&=4\eta+j_0, & H_*&=D\eta+dU_*,\\
  W_*&=1-dU_{*\eta}-2D\eta U_*,
  & Z_*&=-A(1-2\eta U_*)U_*-H_*U_{*\eta}-d\Pi_{0\eta}+4A\eta\Pi_0.
\end{aligned}
\tag{B.1}\label{eq:B.1}
\end{equation}
The function \NSi{NS-F-P144-020}{H_*} has exactly one zero \NSi{NS-F-P144-021}{\eta_0\in(-1,0)}: on
\NSi{NS-F-P144-022}{(-1,1)} the function \NSi{NS-F-P144-023}{H_*/d=D\eta/d+4\eta+j_0} is strictly increasing from
\NSi{NS-F-P144-024}{-\infty} to \NSi{NS-F-P144-025}{+\infty}. Its zero satisfies
\NSi{NS-F-P144-026}{|\eta_0|\asymp j_0}. Moreover
\NSFormula{NS-F-P144-027}{display}{none}
\[
  -W_*=3-8h\eta^2+(1-2h)j_0\eta>2.8,
  \qquad Z_*(\eta_0)\geq cj_0P_*^2>0.
\]
For the second assertion, the term \NSi{NS-F-P144-028}{4A\eta_0\Pi_0(\eta_0)} is positive and bounded below by
\NSi{NS-F-P144-029}{cj_0P_*^2}; the term \NSi{NS-F-P144-030}{-d\Pi_{0\eta}(\eta_0)} is nonnegative, the
\NSi{NS-F-P144-031}{H_*} term vanishes, and the remaining term is \NSi{NS-F-P144-032}{O(j_0)}. The already large
\NSi{NS-F-P144-033}{P_*} absorbs this last term.

Consequently one can first choose \NSi{NS-F-P144-034}{\delta_*>0} so that
\NSi{NS-F-P144-035}{\{|Z_*|\leq\delta_*\}} avoids a neighborhood of \NSi{NS-F-P144-036}{\eta_0}, and then choose
\NSi{NS-F-P144-037}{\sigma_*>0} so small that
\NSFormula{NS-F-P144-038}{display}{B.2}
\begin{equation}
  \chi=\frac{H_*^2}{H_*^2+\sigma_*^2}>.99
  \qquad\text{where }|Z_*|\leq\delta_*.
  \tag{B.2}\label{eq:B.2}
\end{equation}
These choices are made on a slightly larger compact interval \NSi{NS-F-P144-039}{I\supset[-1,1]}. Choose a bounded simply
connected complex neighborhood \NSi{NS-F-P144-040}{\Omega} of \NSi{NS-F-P144-041}{I} on whose closure
\NSi{NS-F-P144-042}{L} and \NSi{NS-F-P144-043}{H_*^2+\sigma_*^2} do not
